\section{Centred groups and the Helling level}\label{sec:centring}

\subsection{Centres of groups}

\begin{lemma}[Fixed sets of compact open subgroups]\label{lem:fix-compact-open}
\slogan{A compact open subgroup of $\AdelicGroup$ has a finite fixed set,
which is a product of local subtrees or midpoints.}
Let $U \subseteq \AdelicGroup$ be a compact open subgroup and let $U_p$ be
the image of $U$ in $\PGL{2}{\pAdicNumbers{p}}$. Then
$\FixedSet{U} = \prod_p \FixedSet{U_p}$, each $\FixedSet{U_p}$ is a nonempty
bounded closed convex subset of $\TreeRealization{p}$ as in
\cref{lem:fixed-shape} (1) or (2), $\FixedSet{U_p} = \{\LocalBaseVertex{p}\}$
for almost all $p$, and $\FixedSet{U}$ is finite.
\end{lemma}

\begin{proof}
The action is coordinatewise, so the fixed set is the product of the
local fixed sets, each of which is nonempty by
\cref{lem:compact-fixes-type}. Since $U$ is open it contains
$\AdelicProjection\{g \in \GL{2}{\ProfiniteIntegers} : g \equiv 1 \bmod M\}$ for
some $M$, so $\FixedSet{U_p}$ is contained in the ball of radius
$\val{p}{M}$ around $\LocalBaseVertex{p}$ by the fixed sets of principal
congruence subgroups (\cref{lem:fixed-principal}).
\end{proof}

\begin{definition}[Centre and Helling level]\label{def:helling-level}
Let $U \subseteq \AdelicGroup$ be a compact open subgroup. Its
\emph{centre} is $\Centre{U} := (\Centre{\FixedSet{U_p}})_p \in \BigPicture$,
with notation as in \cref{lem:fix-compact-open}. Its \emph{Helling level}
$\HellingLevel{U}$ is the product of the primes $p$ for which
$\Centre{U}_p$ is the midpoint of an edge; so $\Centre{U}$ is a point of
type $\HellingLevel{U}$. The group $U$ is \emph{centred} if
$\Centre{U} = \StdConfig{\HellingLevel{U}}$. For a subgroup
$\Gamma \subseteq \RationalProjGroup$ commensurable with $\PSL{2}{\ZZ}$ the
same definitions are applied to the closure of $\lifted{\Gamma}$ in
$\AdelicGroup$ (\cref{lem:closure-commensurable}).
\end{definition}

\begin{lemma}[Centres are equivariant]\label{lem:centre-equivariant}
Let $U \subseteq \AdelicGroup$ be a compact open subgroup and
$g \in \AdelicGroup$. Then $\FixedSet{gUg^{-1}} = g\FixedSet{U}$,
$\Centre{gUg^{-1}} = g\Centre{U}$ and $\HellingLevel{gUg^{-1}} =
\HellingLevel{U}$.
\end{lemma}

\begin{proof}
Equivariance of centres (\cref{lem:centre}); being a midpoint is preserved
by isometries.
\end{proof}

\subsection{Centring}

\begin{proposition}[Centring]\label{prop:centring}
\slogan{Every conjugacy class has centred members, and they form a single
conjugacy class under the Atkin--Lehner group.}
Let $U \subseteq \AdelicGroup$ be a compact open subgroup and
$N := \HellingLevel{U}$. Then there is $g \in \AdelicGroup$ with
$gUg^{-1}$ centred, and $gUg^{-1} \subseteq \MaxCompact{N}$. If $U$ and
$U' = hUh^{-1}$, $h \in \AdelicGroup$, are both centred, then
$h \in \MaxCompact{N}$.
\end{proposition}

\begin{proof}
Choose $g$ with $g\Centre{U} = \StdConfig{N}$ by transitivity on points of
type $N$ (\cref{lem:transitive}); then $gUg^{-1}$ is centred by
equivariance of centres (\cref{lem:centre-equivariant}), and it fixes its
centre. If $U$ and $hUh^{-1}$ are centred then $h\StdConfig{N} =
\StdConfig{N}$.
\end{proof}

\begin{corollary}[Classical centring]\label{cor:classical-centring}
Let $\Gamma \subseteq \RationalProjGroup$ be commensurable with
$\PSL{2}{\ZZ}$ and $N := \HellingLevel{\Gamma}$. Then there is
$\gamma \in \RationalProjGroup$ with $\gamma\Gamma\gamma^{-1}$ centred and
contained in $\GammaZeroStar{N}$. Two centred subgroups of
$\RationalProjGroup$ commensurable with $\PSL{2}{\ZZ}$ are conjugate in
$\RationalProjGroup$ if and only if they are conjugate in
$\GammaZeroStar{N}$.
\end{corollary}

\begin{proof}
As in \cref{prop:centring}, using rational transitivity on points of type
$N$ (\cref{lem:transitive}) and $\MaxCompact{N} \cap (\RationalLiftGroup) =
\GammaZeroStar[lift]{N}$.
\end{proof}

\begin{proposition}[Normalisers are confined]\label{prop:normaliser-confined}
\slogan{The normaliser of a centred group lies in the Atkin--Lehner group
of its Helling level.}
Let $U \subseteq \AdelicGroup$ be a centred compact open subgroup and
$N := \HellingLevel{U}$. Then $\Normalizer{\AdelicGroup}{U} \subseteq
\MaxCompact{N}$. Likewise, for a centred $\Gamma \subseteq \RationalProjGroup$
commensurable with $\PSL{2}{\ZZ}$ one has
$\Normalizer{\RationalProjGroup}{\Gamma} \subseteq \GammaZeroStar{N}$.
\end{proposition}

\begin{proof}
An element of the normaliser preserves $\FixedSet{U}$, hence fixes its
centre $\StdConfig{N}$ by equivariance of centres (\cref{lem:centre}).
\end{proof}

\begin{proposition}[Normalisers of compact open subgroups]
\label{prop:normaliser-compact-open}
\slogan{In $\AdelicGroup$ normalisers of compact open subgroups are compact
open.}
Let $U \subseteq \AdelicGroup$ be a compact open subgroup. Then
$\Normalizer{\AdelicGroup}{U}$ is a compact open subgroup and
$\Normalizer{\AdelicGroup}{U}/U$ is finite.
\end{proposition}

\begin{proof}
Conjugating, we may assume $U$ centred by \cref{prop:centring}. The
normaliser is open since $U$ is, and contained in the compact group
$\MaxCompact{\HellingLevel{U}}$ by \cref{prop:normaliser-confined}.
\end{proof}

\subsection{The Atkin--Lehner image}

\begin{definition}[Atkin--Lehner image]\label{def:al-image}
Let $U \subseteq \AdelicGroup$ be a centred compact open subgroup with
$N := \HellingLevel{U}$. Its \emph{Atkin--Lehner image} is
$\ALimage{U} := \TypeMap(U) \subseteq \TypeGroup{N}$ and
$\ALdim{U} := \dim_{\FF{2}}\ALimage{U}$. The same definitions apply to
centred $\Gamma \subseteq \RationalProjGroup$ commensurable with
$\PSL{2}{\ZZ}$, via $\lifted{\Gamma} \subseteq \GammaZeroStar[lift]{N}$.
\end{definition}

\begin{lemma}[Index formula]\label{lem:index-formula}
Let $U \subseteq \AdelicGroup$ be a centred compact open subgroup with
$N := \HellingLevel{U}$. Then
$\Index{\MaxCompact{N}}{U} = 2^{\NumPrimeFactors{N}-\ALdim{U}}
\Index{\GammaZero[ad]{N}}{U \cap \GammaZero[ad]{N}}$.
\end{lemma}

\begin{proof}
The type map identifies $\MaxCompact{N}/\GammaZero[ad]{N}$ with
$\TypeGroup{N}$ (\cref{lem:U_N}), and $U\cap\GammaZero[ad]{N}$ is the kernel
of $\TypeMap|_{U}$.
\end{proof}

\begin{remark}
The Atkin--Lehner image need not be all of $\TypeGroup{N}$: a midpoint
centre can come from an inverted edge, but also from a fixed subtree of
odd diameter. For example $\GammaZero[ad]{p}$ is centred with Helling level
$p$ and $\ALdim{\GammaZero[ad]{p}} = 0$ (\cref{ex:centring-examples}).
\end{remark}

\begin{example}[Centring examples]\label{ex:centring-examples}
Let $p$ be a prime.
\begin{enumerate}
\item $\FixedSet{\GammaZero[ad]{p}}$ is the edge from $\BaseVertex$ to
  $\StdVertex{p}$ (\cref{lem:local-fixed-points}), so $\GammaZero[ad]{p}$
  is centred, $\HellingLevel{\GammaZero[ad]{p}} = p$ and
  $\ALimage{\GammaZero[ad]{p}} = 0$. The group $\GammaZeroStar[ad]{p}$ is
  centred with the same Helling level and $\ALdim{} = 1$.
\item Let $p \geq 3$ and $k \geq 1$. Then $\FixedSet{\GammaZero[ad]{p^{k}}}$
  is the local snake from $\LocalBaseVertex{p}$ to $\LocalStdVertex{p^{k}}{p}$:
  at an interior vertex $\LocalStdVertex{p^{j}}{p}$ the group acts on the
  remaining $p-1$ neighbours through the diagonal torus of
  $\GL{2}{\FF{p}}$, which has no fixed point there when $p \geq 3$, and at
  the endpoints one argues as for the Iwahori subgroup. Hence
  $\HellingLevel{\GammaZero[ad]{p^{k}}} = 1$ for $k$ even and $p$ for $k$
  odd.
\item Let $p \geq 3$. Conjugating $\GammaZero[ad]{p^{2}}$ by
  $\Diagonal{1}{p^{-1}}$, which maps $\StdVertex{p}$ to $\BaseVertex$, gives
  the centred group
  $\{\smallMat{a}{b}{c}{d} : b \equiv c \equiv 0 \bmod p\} =
  \GammaZero[ad]{p}\cap\GammaUpperZero[ad]{p}$, where $\GammaUpperZero[ad]{p}$
  is defined like $\GammaZero[ad]{p}$ with the condition on $b$ instead of
  $c$. It has Helling level $1$.
\item Let $p \geq 3$. Conjugating $\GammaZero[ad]{p^{3}}$ by
  $\Diagonal{1}{p^{-1}}$, which
  maps $\StdVertex{p}, \StdVertex{p^{2}}$ to $\BaseVertex, \StdVertex{p}$,
  gives the centred group $\{\smallMat{a}{b}{c}{d} : b \equiv 0 \bmod p,\
  c \equiv 0 \bmod p^{2}\}$, of Helling level $p$.
\end{enumerate}
\todo[inline]{For $p \in \{2, 3\}$ the fixed sets of $\GammaZero{p^{k}}$
can be larger than the snake (e.g.\ $\GammaZero{4}$ is conjugate to
$\GammaFull{2}$, and classically $\GammaZero{9}$ to $\GammaFull{3}$); this
is the source of the divisors of $24$ in Conway's normaliser theorem. Work
out the centres in these cases.}
\end{example}

\begin{remark}[Moduli interpretation]\label{rem:moduli-up-to-isogeny}
The group $\PGL{2}{\FiniteAdeles{\QQ}}$ acts on the projective level
structures of an isogeny class of elliptic curves, and the vertices of
$\BigPicture$ attached to an isogeny class are the curves in it (up to
isomorphism). From this point of view a point of
$\OpenModCurve{\MaxCompact{N}}$ is an unordered pair of $N$-isogenous curves,
up to the dual isogeny; the rational points of $\ModCurve{\GammaZeroStar{N}}$
are, away from cusps and CM points, essentially $\QQ$-curves of degree $N$
\cite{Elkies2004}. For compact open subgroups of nontrivial type there is
no universal elliptic curve.
\end{remark}
